% Post-checkpoint 07.  Independent audit requested before integration.
\begin{proposition}[Proved: a uniform first-cohomology bound does not imply finite generation]
\label{D:uniform-H1-obstruction}
The group $G$ of Example~\ref{D:pronilpotent-obstruction} satisfies
\[
 \dim_{\mathbb F}H^1_{
 \mathrm{cont}}(G,M)\leq\dim_{\mathbb F}M
\]
for every finite field $\mathbb F$ and every finite-dimensional
continuous $\mathbb F[G]$-module $M$.  Nevertheless $G$ is not
topologically finitely generated.  Moreover, for every open subgroup
$H\subset G$ and every finite continuous $\mathbb F[H]$-module $V$,
\[
 \dim_{\mathbb F}H^1_{
 \mathrm{cont}}(H,V)\leq [G:H]\dim_{\mathbb F}V.
\]
Thus even these estimates together with smallness and the previously
proved restrictions on all open subgroups do not settle the
finite-generation question.
\end{proposition}
\begin{proof}
We first give a bound for a finite direct product
$P=\prod_{i=1}^r S_i$ of perfect two-generated finite groups.
For any field $F$ and any finite-dimensional $F[P]$-module $M$,
\begin{equation}\label{D:perfect-product-H1}
 \dim_F H^1(P,M)\leq\dim_F M.
\end{equation}
We may extend scalars to an algebraic closure: kernels and images
in the finite group cochain complex commute with flat scalar
extension.  A composition series and the long exact cohomology
sequence reduce the inequality to irreducible modules.

Over an algebraically closed field every irreducible module for
a finite direct product has the form $\bigotimes_i M_i$, where
each $M_i$ is an irreducible $S_i$-module.  This assertion also
holds in modular characteristic.  Indeed, the ideals obtained
by tensoring the Jacobson radical of any factor group algebra
with the remaining algebras are nilpotent, and hence annihilate
an irreducible module.  After quotienting by them, the assertion
reduces to the tensor product of split matrix algebras.

If all factors act trivially, perfectness gives $H^1(P,M)=0$.
If exactly one factor $S$ acts nontrivially, the other factors
have no homomorphism to the additive group of $M$.  Restriction
therefore identifies $H^1(P,M)$ with $H^1(S,M)$.
A cocycle on $S$ is determined by its values on two generators,
so $\dim Z^1(S,M)\leq2\dim M$.  Since $M^S=0$,
$\dim B^1(S,M)=\dim M$, proving the required inequality.

Suppose that at least two factors act nontrivially.  Write
$P=S\times T$ and $M=V\otimes W$, where $V$ is a nontrivial
irreducible $S$-module and $W$ is a nontrivial irreducible
$T$-module.  Restriction gives an injection
\[
 H^1(P,M)\hookrightarrow H^1(S,M)^T
 =\bigl(H^1(S,V)\otimes W\bigr)^T=0.
\]
For clarity, the injection follows without any semisimplicity
assumption: a cocycle whose restriction to $S$ is a coboundary
can be adjusted to vanish on $S$, after which its values on
$T$ lie in $M^S=0$.  The image is $T$-invariant because the
cocycle identity on commuting elements makes conjugation differ
by a coboundary.  The last equality follows from $W^T=0$.
This proves~\eqref{D:perfect-product-H1}.

The alternating groups $A_n$ are perfect for $n\geq5$ and
two-generated.  Explicit generators are
\[
 \begin{cases}
  (1\,2\,\cdots\,n),\ (1\,2\,3),& n\ \text{odd},\\
  (2\,3\,\cdots\,n),\ (1\,2\,3),& n\ \text{even}.
 \end{cases}
\]
For odd $n$, conjugation gives the consecutive three-cycles,
which generate $A_n$ by induction.  For even $n$, conjugation
gives $(1\,i\,i+1)$; the products of consecutive such cycles
give $(i\,i+1\,i+2)$ on $\{2,\ldots,n\}$.  These generate
$A_{n-1}$, and the remaining cycle joins the letter $1$.
The induction starts with $A_3$, and $A_4$ follows from
$(1\,2\,3),(2\,3\,4)$: their generated group is transitive
and has a point stabilizer of order at least three, hence has
order twelve.

Every finite continuous $G$-module factors through a finite
subproduct $P$, whose complementary product $K$ acts trivially.
Every continuous homomorphism from $K$ to a finite abelian
group vanishes: it vanishes on each perfect factor and then on
their dense finite-support subgroup.  Consequently restriction
and inflation identify $H^1(G,M)$ with $H^1(P,M)$, and
\eqref{D:perfect-product-H1} proves the first assertion.
The failure of finite generation and all the smallness and
pronilpotent assertions were proved in
Example~\ref{D:pronilpotent-obstruction}.  Finally Shapiro's
lemma \cite[II, Proposition 1.11, p.~62, and Proposition 4.2,
p.~87]{MilneCFT2020D}, applied to the finite induced module, gives
\[
 H^1(H,V)=H^1(G,\operatorname{Ind}_H^G V),\qquad
 \dim\operatorname{Ind}_H^G V=[G:H]\dim V,
\]
proving the open-subgroup estimate.

The exact certificate
\nolinkurl{parts/D/code/check_uniform_h1_obstruction.py}
checks the displayed generators for $5\leq n\leq8$, the
four-dimensional deleted permutation module of $A_5$ over
$\mathbb F_2$, and its outer tensor square for $A_5^2$.
It gives $\dim H^1=0$ in both module examples, checks the
first calculation independently with all $14\,400$ bar-cocycle
equations, and verifies $2\,803\,802$ exact assertions.
The mathematical proof above, rather than these finite checks,
establishes the claim for every module and the infinite product.
This is an abstract obstruction to a proof strategy, not a
geometric counterexample to D.1 or D.2(ii).
\end{proof}

\begin{proposition}[Proved: the two uniform bounds needed for profinite generation]
\label{D:cohomology-crown-criterion}
Let $G$ be a profinite group.  For a finite quotient $F$ and a
non-Frattini chief factor $A$ of $F$, let $\delta_F(A)$ be the
number of non-Frattini chief factors in its $F$-equivalence
class.  Equivalently, this is the multiplicity in the crown
quotient $(L_A)_{\delta_F(A)}$; here $L_A$ is the associated
monolithic primitive group and
\[
 (L_A)_m=\{(l_1,\ldots,l_m)\in L_A^m:
                 l_1A=\cdots=l_mA\}.
\]
Then $G$ is topologically finitely generated if and only if
there is an integer $c\geq1$ such that both conditions hold:
\begin{enumerate}
\item for every prime $\ell$ and every finite irreducible
continuous $\mathbb F_\ell[G]$-module $V$,
\[
 \dim_{\mathbb F_\ell}H^1(G,V)
       \leq c\dim_{\mathbb F_\ell}V;
\]
\item for every finite quotient $F$ and every nonabelian
chief factor $A$ of $F$,
\[
 \delta_F(A)\leq |A|^c.
\]
\end{enumerate}
These conditions imply that $G$ is generated by at most $c+2$
elements.  No assertion that either uniform condition is already
known geometrically in the residual D.1/D.2(ii) cases is made.
\end{proposition}
\begin{proof}
We use the finite crown reduction of
\cite[Lemma 2.5 and Proposition 2.6, pp.~10--13]{LucchiniMarionTracey2020D}.
The required inputs are these.  A finite noncyclic group minimal
among its quotients with a given generator number has the form
$L_k$, where $L$ is monolithic primitive with socle $A$.
If $A=S^n$ is nonabelian, then
\[
 d(L)=\max\{2,d(L/A)\},\qquad
 k\leq\psi_L(s)\Longrightarrow d(L_k)\leq s
 \quad(s\geq d(L)),
\]
where
\[
 \psi_L(s)\geq
 \frac{|A|^{s-1}}{2n|\operatorname{Out}(S)|},
 \qquad |\operatorname{Out}(S)|\leq\log_2|S|.
\]
The crown quotient and multiplicity convention are those of
Proposition 2.4, p.~10; the two inequalities are equations
(2.1)--(2.2), p.~12, and the outer-automorphism bound is stated
on p.~13.  These finite-group inputs include the classification
of finite simple groups.  The following deduction is explicit.

If $G$ is $d$-generated, cocycles are determined by their values
on $d$ generators, proving the first bound with $c\geq d$.
Every crown quotient $L_m$ is also $d$-generated.  The $d$-tuple
in each coordinate has the same prescribed image in $(L/A)^d$,
so there are at most $|A|^d$ possible coordinate tuples.
Two identical tuples would force the generated group into a
proper equality diagonal.  Therefore $m\leq|A|^d$, proving
the second bound.

Conversely, fix a finite quotient $F$ and put $d=d(F)$.
There is nothing to prove if $d\leq2$.  Choose a quotient $Q$
of smallest order with $d(Q)=d$.  Every proper quotient of $Q$
has fewer than $d$ generators.  The crown reduction gives
$Q=L_k$, with socle $A^k$ and $H=L/A$; in particular
$d(H)<d$.  The factor $A$ occurs among the chief factors of
$F$ and $k\leq\delta_F(A)$.

First suppose that $A$ is abelian.  Thus $L=A\rtimes H$, and
$A$ is a faithful irreducible $\mathbb F_\ell[H]$-module.
Write $D=\operatorname{End}_H(A)$, $r=\dim_D A$ and
$s=\dim_D H^1(H,A)$.  Since $Q=A^k\rtimes H$ is split,
restriction of cocycles gives
\[
 \dim_D H^1(Q,A)=s+k.
\]
Indeed, every $H$-equivariant homomorphism $A^k\to A$
extends to a cocycle on the semidirect product, and the kernel
is the inflated $H^1(H,A)$.  Inflation from $Q$ to $G$ is
injective.  Condition (1), with dimensions converted from
$\mathbb F_\ell$ to $D$, therefore gives $s+k\leq cr$.

We record directly why this bounds generators, to avoid any
unstated restriction on $d(H)$.  Choose an integer
$b\geq d(H)$ and generators $h_1,\ldots,h_b$ of $H$,
allowing repetitions of the identity.  Evaluating cocycles
on this tuple embeds $Z^1(H,A)$ in $A^b$.  Its $D$-dimension
is $s+\theta r$, where $\theta=0$ if $H=1$ and $\theta=1$
otherwise.  If
\[
 k\leq br-s-\theta r,
\]
choose $k$ independent vectors in $A^b/Z^1(H,A)$ and use
their representatives as the coordinate values of lifts
$(a_i,h_i)\in A^k\rtimes H$.  These lifts generate $Q$.
Otherwise their intersection with $A^k$ is a proper
$H$-submodule, and some nonzero $H$-map $A^k\to A$
annihilates that intersection.  Their image in $A\rtimes H$
would then be a complement, hence the graph of a cocycle,
contradicting the chosen independence.
Taking $b=\max\{d(H),c+1\}$ meets the inequality.
Since $d(H)<d$, it follows that $d\leq c+1$.

If $A=S^n$ is nonabelian, we have
$d(L)=\max\{2,d(H)\}\leq d-1$.
Suppose that $d\geq c+3$ and take $s=d-1$ in the displayed
crown bound.  Since
\[
 2n|\operatorname{Out}(S)|
 \leq2\log_2|A|\leq |A|,
\]
we obtain
\[
 \psi_L(d-1)\geq |A|^{d-3}
      \geq |A|^c\geq k.
\]
This forces $d(Q)\leq d-1$, a contradiction.  Hence
$d(F)\leq c+2$ in every case.

Finally, in the compact space $G^{c+2}$ the sets of tuples
generating each finite quotient are nonempty clopen sets
with the finite-intersection property.  A tuple in their
intersection topologically generates $G$.
\end{proof}

\begin{remark}[The exact remaining cohomological route]
For the groups in D.1 and D.2(ii), our smallness proofs give
finiteness for each fixed finite torsor group and each fixed
finite module.  They provide neither uniform constant in
Proposition~\ref{D:cohomology-crown-criterion}.  In the
codimension-two case, $H^1(\pi_1(U),V)$ is the ordinary
\'etale $H^1(U,\mathcal V)$.  The degree-one Leray sequence
for $j:U\hookrightarrow X$ separates a global term
$H^1(X,j_*\mathcal V)$ from a local term in
$H^0(X,R^1j_*\mathcal V)$.  Constructibility and properness
give pointwise finiteness, without a proved estimate linear
in the rank as $\mathcal V$ varies.  For D.1 one must retain
the fixed-boundary tame cohomology classes: ordinary
$H^1(U,\mathbb F_p)$ can already be infinite on $\mathbb A^1$.

Even a successful uniform linear cohomology estimate would
leave condition (2).  In the preceding alternating-product
example, the finite quotient $A_n^{m_n}$ has
$\delta(A_n)=m_n=|A_n|^n+1$, so that condition fails for
every fixed exponent while condition (1) holds with $c=1$.
Thus a full argument needs a geometric bound on these
nonabelian multiplicities, or a different geometric reduction
that excludes them.  We do not infer such a bound from
Gabber's constructibility theorem alone.
\end{remark}
